\begin{table}[htbp]
    \centering
    \begin{tabular}{@{}lrrr@{}}
        \toprule
        Spectral curve                           & $n_{\mathrm{max}}$ & $N_{\mathrm{res}}$ & $N_{\mathrm{cor}}$ \\
        \midrule
        Generic cubic two-matrix                 & 4                  & 10                 & 6 \\
        Quartic two-matrix $(g,g,g)$, $(g,-g,g)$ & 6                  & 11                 & 1 \\
        Quartic two-matrix $(g,g,-g)$            & 7                  & 27                 & 4 \\
        Quartic two-matrix $(-g,g,g)$            & 7                  & 27                 & 4 \\
        Cubic three-matrix                       & 5                  & 64                 & 7 \\
        3-state Potts, fixed                     & 5                  & 64                 & 7 \\
        3-state Potts, mixed                     & 6                  & 245                & 17 \\
        3-state Potts, free                      & 6                  & 245                & 36 \\
        General 3-state Potts, fixed             & 6                  & 633                & 84 \\
        \bottomrule
    \end{tabular}
    \caption{Levels and equation counts in the accompanying notebook derivations.
        Here $n_{\mathrm{max}}$ is the largest resolvent-word level in the generated
        system, and $N_{\mathrm{res}}$ and $N_{\mathrm{cor}}$ count the nonzero
        resolvent and correlator equations after the notebook's symmetry identifications
        and removal of duplicates, before elimination. The two quartic cases in the
        second row each have the stated counts; their signed-sector zero rules have
        already been applied. These are sufficient systems used in the notebooks,
        not claims of minimal equation counts or minimal levels.}
    \label{tab:curve-counts}
\end{table}

Table~\ref{tab:curve-counts} records the size of the loop systems used below. For cubic potentials
the maximal variation-word length is $n_{\mathrm{max}}-2$, and for quartic potentials it is
$n_{\mathrm{max}}-3$. The generic cubic two-matrix notebook generates ten resolvent equations; the
eight equations displayed below form a smaller sufficient subset.
